\documentclass[11pt]{article}
\pdftrailerid{}
\usepackage[T1]{fontenc}
\usepackage{lmodern}
\ifdefined\pdfglyphtounicode
  \pdfglyphtounicode{intersectiondisplay}{22C2}
  \pdfglyphtounicode{intersectiontext}{22C2}
  \pdfglyphtounicode{parenleftBigg}{0028}
  \pdfglyphtounicode{parenleftbig}{0028}
  \pdfglyphtounicode{parenrightBigg}{0029}
  \pdfglyphtounicode{parenrightbig}{0029}
  \pdfglyphtounicode{summationdisplay}{2211}
\fi
\usepackage[margin=1in]{geometry}
\usepackage{amsmath,amssymb,amsthm,mathtools}
\usepackage{aliascnt}
\usepackage{microtype}
\usepackage{needspace}
\usepackage{tikz}
\usetikzlibrary{arrows.meta,positioning}
\usepackage{enumitem}
\usepackage{xcolor}
\usepackage[colorlinks=true,linkcolor=blue!45!black,citecolor=blue!45!black,
  urlcolor=blue!45!black,pdftitle={Bounded-degree coboundary expanders in every dimension},
  pdfauthor={OpenAI}]{hyperref}
\usepackage[nameinlink,noabbrev]{cleveref}

\newtheorem{theorem}{Theorem}[section]
\newaliascnt{proposition}{theorem}
\newtheorem{proposition}[proposition]{Proposition}
\aliascntresetthe{proposition}
\newaliascnt{lemma}{theorem}
\newtheorem{lemma}[lemma]{Lemma}
\aliascntresetthe{lemma}
\newaliascnt{corollary}{theorem}
\newtheorem{corollary}[corollary]{Corollary}
\aliascntresetthe{corollary}
\crefname{theorem}{Theorem}{Theorems}
\crefname{proposition}{Proposition}{Propositions}
\crefname{lemma}{Lemma}{Lemmas}
\crefname{corollary}{Corollary}{Corollaries}
\newcommand{\F}{\mathbb F_2}
\newcommand{\Z}{\mathbb Z}
\newcommand{\GL}{\operatorname{GL}}
\newcommand{\SL}{\operatorname{SL}}
\newcommand{\CC}{\operatorname{CC}}
\newcommand{\sr}{\operatorname{sr}}
\newcommand{\Tor}{\operatorname{Tor}}
\newcommand{\im}{\operatorname{im}}
\newcommand{\dist}{\operatorname{dist}}
\newcommand{\skel}{\operatorname{skel}}
\newcommand{\Stab}{\operatorname{Stab}}
\newcommand{\wtH}{\widetilde H}
\newcommand{\wtC}{\widetilde C}
\newcommand{\eps}{\varepsilon}
\DeclarePairedDelimiter{\norm}{\lVert}{\rVert}
\numberwithin{equation}{section}
\setlist[enumerate]{leftmargin=*,label=\textup{(\roman*)}}
\setlength{\emergencystretch}{1.5em}


\title{Bounded-degree coboundary expanders in every dimension}
\author{OpenAI}
\date{September 24, 2026}
\begin{document}
\maketitle
\begin{abstract}
For every integer $d\ge3$, we construct arbitrarily large finite
$d$-dimensional simplicial complexes with uniformly bounded vertex degrees
and uniform $\mathbb F_2$ coboundary expansion in every degree below $d$.
Together with the known graph and two-dimensional cases, this establishes
the existence of bounded-degree $\mathbb F_2$ coboundary expanders in every
positive dimension.
\end{abstract}
\section{Introduction}
\label{sec:introduction}

In an expander graph, the edge boundary of a set of vertices controls
its distance from the empty set or the whole vertex set. Coboundary
expansion extends this condition to sets of faces in a simplicial
complex. Over $\F$, the coboundary of a set of $i$-faces consists of
the $(i+1)$-faces containing an odd number of them. We construct
complexes of increasing size for which this boundary controls distance
from the space of coboundaries, uniformly in the size and with a bounded
number of top-dimensional faces at each vertex.

Let $Y$ be a finite pure $s$-dimensional simplicial complex, and let $Y(i)$
denote its faces with $i+1$ vertices. Purity means that every face is
contained in an $s$-face. Write $C^i(Y;\F)$ for the functions from
$Y(i)$ to $\F$. We use the weights
\begin{equation}\label{eq:weights}
 w_i^Y(\sigma)=
 \frac{\#\{\tau\in Y(s):\sigma\subseteq\tau\}}
 {\binom{s+1}{i+1}\,|Y(s)|},
 \qquad
 \norm{f}_{Y,i}=\sum_{\sigma:\,f(\sigma)\ne0}w_i^Y(\sigma).
\end{equation}
Equivalently, sample an $s$-face uniformly and then one of its $i$-faces
uniformly; $\norm{f}_{Y,i}$ is the probability that $f$ is nonzero there.
All simplicial chains and cochains have coefficients in $\F$ unless
specified otherwise. The coboundary is
\[
 (\delta_i f)(\tau)=\sum_{\substack{\sigma\subset\tau\\|\sigma|=i+1}}
 f(\sigma)\quad\text{in }\F.
\]
Write $Z^i(Y)=\ker\delta_i$ and
$B^i(Y)=\im\delta_{i-1}$, taking $B^0(Y)$ to be the constant functions.
For a subspace $A\subseteq C^i(Y;\F)$, put
\[
 \dist_{Y,i}(f,A)=\min_{a\in A}\norm{f-a}_{Y,i}.
\]
Thus $Z^0(Y)=B^0(Y)$ when $Y$ is connected.

\begin{theorem}\label{thm:main}
For every integer $d\ge3$, there exist constants $D<\infty$, $\eps>0$,
and finite connected pure $d$-dimensional simplicial complexes $X_m$,
with $|X_m(0)|\longrightarrow\infty$, such that every vertex belongs to
at most $D$ faces of dimension $d$ and, for every
$f\in C^i(X_m;\F)$,
\begin{equation}\label{eq:goal}
 \norm{\delta_i f}_{X_m,i+1}
 \ge \eps\,\dist_{X_m,i}(f,B^i(X_m))
 \qquad(0\le i<d).
\end{equation}
The constants are independent of $m$, $i$, and $f$.
\end{theorem}

The distinction between cocycles and coboundaries is essential. If
$f\in Z^i(Y)\setminus B^i(Y)$, its coboundary vanishes while its distance
from $B^i(Y)$ is positive. Thus \eqref{eq:goal} requires the vanishing of
$H^i(Y;\F)=Z^i(Y)/B^i(Y)$, with reduced cohomology understood in degree
zero. A uniform inequality against $Z^i(Y)$ alone allows nonzero
cohomology. Our proof supplies both that quantitative inequality and
the required vanishing for the same complexes.

\paragraph{Prior work.}
The higher-dimensional coisoperimetric viewpoint developed through the
random-complex work of Linial and Meshulam \cite{LM2006} and Meshulam and
Wallach \cite{MW2009}. Their models retain a complete lower skeleton and
study the disappearance of cohomology as top-dimensional faces are
added. Gromov connected cohomological filling inequalities with
topological overlap and asked for arbitrarily large complexes with
bounded vertex incidence and uniform filling bounds
\cite[Sections~2.3--2.5 and~2.14]{Gromov2010}. The full coboundary
expansion problem also demands vanishing cohomology; a recent formulation
is \cite[Problem~1.3]{CL2025}.

There are two different degree conditions in this history. Lubotzky and
Meshulam constructed random Latin-square coboundary expanders in
dimension two \cite{LM2015}; Lubotzky, Luria, and Rosenthal constructed
random Steiner-system coboundary expanders in every dimension
\cite{LLR2019}. These results bound the number of top-dimensional faces
through a codimension-one face. The Latin-square complexes have a
complete tripartite $1$-skeleton, and the Steiner-system complexes have
a complete $(d-1)$-skeleton. In dimensions at least two, vertex degrees
grow with the size. The degree bound in \cref{thm:main} is at vertices.

For bounded vertex degree, Kaufman, Kazhdan, and Lubotzky established
cosystolic and topological expansion for two-dimensional skeleta of
Ramanujan complexes \cite{KKL2016}. Evra and Kaufman obtained
bounded-degree cosystolic expanders in every dimension through a
local-to-global theorem \cite{EK2024}. The associated expansion methods
include the building-like complexes of Lubotzky, Meshulam, and Mozes
\cite{LMM2016}, the strongly symmetric complexes of Kaufman and Oppenheim
\cite{KO2021}, the coefficient-independent methods of Kaufman and Mass
\cite{KM2021}, and the quantitative local-to-global results of Dikstein
and Dinur \cite{DD2024}. We use only the distance-to-cocycles part of
cosystolic expansion, stated precisely in \eqref{eq:cosystolic}.

Chapman and Lubotzky proved the existence of bounded-degree
two-dimensional coboundary expanders over $\F$
\cite[Theorem~1.9]{CL2025}. In higher ambient dimensions, Kaufman,
Oppenheim, and Weinberger \cite{KOW2025}, and Oppenheim and
Valentiner-Branth \cite{OVBNonAbelian2025}, proved degree-one coboundary
expansion for coset complexes with suitable non-Abelian coefficients.
Those degree-one results and the cosystolic constructions above do not
give the vanishing in every degree required by \cref{thm:main}.
Together with the two-dimensional result and the graph case, our theorem
gives bounded-degree $\F$ coboundary expanders in every positive
dimension.

Our complexes belong to an established algebraic construction program.
Kaufman and Oppenheim developed locally spectral expanding coset
geometries \cite{KO2023}; O'Donnell and Pratt used Chevalley groups
\cite{OP2022}; and Grave de Peralta and Valentiner-Branth developed
Kac--Moody--Steinberg group constructions \cite{GdPVB2025}.
Oppenheim and Valentiner-Branth proved the cosystolic expansion theorem
for these KMS complexes that supplies our quantitative input \cite{OVB}.
Their affine-type-$A$ polynomial matrix construction
\cite[Example~4.13]{OVB} already uses the simple root groups appearing
here, with quotients by irreducible polynomials. Powers of $t$ also occur,
with different local groups, in the degree-one construction of
\cite[Section~2.5 and Theorem~6.3]{KOW2025}.
For the local groups used here, we pass to the tower of quotients modulo
$t^m$ and prove the cohomology vanishing needed to obtain coboundary
expansion in all degrees below $d$.

\paragraph{Proof strategy.}
Fix $d$ and put $n=2d$ and $N=2d+1$. Over a sufficiently large fixed
odd prime field $k$, let $U\le\SL_N(k[t])$ be the inverse image of the
upper unitriangular group under evaluation at zero. We construct an
infinite $n$-dimensional coset complex $L$ for $U$ and its finite
congruence quotients $T_m$. The desired $X_m$ is the $d$-skeleton of
$T_m$. The local groups are fixed finite groups, which gives uniform
degree bounds; the congruence quotients grow with $m$.

The first step is geometric. We identify $L$ with the subcomplex
opposite a chamber in an affine twin building. A filtration by
codistance expresses the relative chain groups of each attachment in
terms of a spherical opposition complex. The opposition-complex input
of \cite{OVB} makes these attachments acyclic in the required range.
Since the whole affine building is contractible, $L$ has vanishing
reduced homology below degree $n-1$. Its simplex stabilizers have odd
order, so its chain complex gives a partial projective resolution over
$\F[U]$. The orbit complex $U\backslash L$ is a single simplex, which
yields the corresponding homology vanishing for $U$.

The second step passes from $U$ to each congruence subgroup. Morrow's
pro Tor-unitality theorem \cite{Morrow} and Iwasa's homology pro-stability
theorem \cite{Iwasa} apply to the ideals $t^m k[t]$. Their initial
conclusion concerns an inverse system, whereas we need a statement at
every level. Odd-index transfer makes the homology maps towards
shallower levels surjective over $\F$. The cokernels of stabilization
are therefore both pro-zero and connected by surjective transition
maps, so each cokernel vanishes. Two stabilization steps then show that
ambient conjugation acts trivially in degrees $1,\ldots,d-1$.
Averaging over the finite odd-order quotient transfers the vanishing
from $U$ to the congruence subgroup, and hence to $T_m$.

Finally, the KMS theorem gives a uniform inequality against cocycles
on the $(n-1)$-skeleton of $T_m$. Its weighted norms agree with the
corresponding norms on $T_m$ because every panel lies in the same
number of chambers. Passing to the $d$-skeleton requires a further
comparison: the number of chambers through a $d$-face is bounded above
and below, and the normalization contributes the same factor in two
successive cochain degrees. This factor cancels in the expansion
inequality. The cohomology vanishing replaces cocycles by coboundaries
and completes the proof.

\paragraph{Organization.}
Section~\ref{sec:construction} defines the congruence complexes.
Section~\ref{sec:geometry} proves the local group identities and the
acyclicity of $L$. Section~\ref{sec:homology} proves the cohomology
vanishing of the finite quotients. Section~\ref{sec:expansion} checks
the KMS hypotheses, compares skeleton weights, and proves
\cref{thm:main}.

\section{The congruence complexes}
\label{sec:construction}

We now define the complexes appearing in \cref{thm:main}. Fix
\[
 n=2d,\qquad N=n+1=2d+1,
\]
and choose a sufficiently large odd prime $q$, depending only on $N$.
The precise largeness requirements will be stated where they are used.
Set $k=\mathbb F_q$, $R=k[t]$, and
\[
 \Gamma_a(m)=\ker\bigl(\SL_a(R)\longrightarrow\SL_a(R/(t^m))\bigr).
\]
Let $U\le\SL_N(R)$ be the inverse image, under evaluation at $t=0$, of
the upper unitriangular subgroup of $\SL_N(k)$. Put
\[
 S=\{0,1,\ldots,N-1\},\qquad
 x_i=e_{i,i+1}(1)\ (1\le i<N),\qquad x_0=e_{N,1}(t),
\]
where $e_{ab}(f)=I+fE_{ab}$, and define
\begin{equation}\label{eq:Hi}
 H_i=\langle x_j:j\in S\setminus\{i\}\rangle.
\end{equation}

For a group $G$ and subgroups $(K_i)_{i\in S}$, the \emph{coset complex}
$\CC(G;(K_i))$ has vertices the tagged left cosets $gK_i$; a collection
of vertices is a face when the corresponding cosets have nonempty common
intersection. Its chambers are the sets $\{gK_i:i\in S\}$, so it is pure
of dimension $|S|-1$. For $m\ge2$, set
\begin{equation}\label{eq:construction}
 G_m=U/\Gamma_N(m),\qquad
 T_m=\CC(G_m;(\overline H_i)_{i\in S}),\qquad
 X_m=\skel_d T_m.
\end{equation}
The bar denotes reduction modulo $t^m$. Here $\Gamma_N(m)\subseteq U$,
and $G_m$ is the finite image of $U$ in $\SL_N(R/(t^m))$.


\section{The infinite coset complex}
\label{sec:geometry}

Put $L=\CC(U;(H_i)_{i\in S})$. We will prove
\begin{equation}\label{eq:L-acyclic}
 \wtH_j(L;\F)=0\qquad(0\le j<n-1).
\end{equation}
This is the acyclicity needed to compare group homology with the homology
of the finite complexes. We first identify the finite stabilizers, then
realize $L$ inside an affine building and compute the relative chains of
a filtration. Connectedness will also show that the $x_i$ generate $U$.

\subsection{The finite local groups}

Regard $S$ as the edges of the directed cycle
$1\to2\to\cdots\to N\to1$, with edge $0$ equal to $N\to1$.
For $J\subsetneq S$, put $U_J=\langle x_j:j\in J\rangle$ and
$U_{\varnothing}=\{1\}$. Between distinct indices $a,b$ there is a unique
simple directed path in the cycle. Set
\[
 V_{ab}(J)=
 \begin{cases}
 k,&\text{if that path lies in $J$ and avoids edge $0$},\\
 kt,&\text{if that path lies in $J$ and uses edge $0$},\\
 0,&\text{otherwise}.
 \end{cases}
\]

\Needspace{12\baselineskip}
\begin{lemma}\label{lem:patterns}
For every proper $J\subsetneq S$,
\begin{equation}\label{eq:pattern}
 U_J=\{g:g_{aa}=1,\ g_{ab}\in V_{ab}(J)\text{ for }a\ne b\}.
\end{equation}
These are finite groups, and
\begin{equation}\label{eq:intersections}
 U_J\cap U_K=U_{J\cap K},\qquad
 \bigcap_{i\in A}H_i=U_{S\setminus A}\quad(\varnothing\ne A\subseteq S).
\end{equation}
For $m\ge2$, reduction modulo $t^m$ is injective on all $U_J$ and preserves
all these intersections. In particular,
\begin{equation}\label{eq:M}
 |H_i|=M:=q^{\binom N2}\quad(i\in S),\qquad |U_{\{i\}}|=q.
\end{equation}
\end{lemma}

\begin{proof}
Deleting an edge outside $J$ cuts the cycle into a linear order. In this
order the right side of \eqref{eq:pattern} is a product of unitriangular
groups, one for each path component of $J$. Factors of $t$ correspond to
conjugation by a diagonal matrix over $k[t,t^{-1}]$. In particular the
displayed matrices form a group.

Because $q$ is prime, powers of $x_j$ provide every coefficient in its
simple root group. Iterated commutators
\[
 [e_{ab}(f),e_{bc}(g)]=e_{ac}(fg)
 \qquad(a,b,c\text{ distinct})
\]
provide every allowed path entry. Unitriangular elimination in the cut
order then gives every matrix on the right of \eqref{eq:pattern}.
This proves the equality. The path descriptions also show the
intersection identities. Finally, every entry has degree at most one;
reduction modulo $t^m$, for $m\ge2$, neither identifies distinct such
entries nor changes their allowed subspaces. The same intersection
identities therefore hold after reduction.
\end{proof}

\subsection{The twin-building model}

The matrix calculation identifies the stabilizers we need. To obtain
acyclicity, we now place the coset complex in the classical affine
twin building of $\SL_N(k[t,t^{-1}])$. We use the lattice construction
of Abramenko--Van Maldeghem
\cite[Section~4, Proposition~4.4 and Remark~4.6]{AVM2001};
see also \cite[Lemma~5]{Abramenko} and
\cite{AbramenkoBrown}. Its two halves are the
affine buildings at $0$ and at infinity, both of type $\widetilde A_n$
and dimension $n$. Their chambers have one vertex of each type in $S$.
The stabilizers of opposite base chambers are
\[
 B_+=\{g\in\SL_N(k[t]):g(0)\text{ is upper triangular}\},
\]
and the analogous group $B_-$ over $k[t^{-1}]$ with lower triangular
reduction at infinity. They satisfy $B_+=U\rtimes T$ and
$B_+\cap B_-=T$, where $T$ is the constant diagonal torus.

Fix the positive base chamber $c_+$. A chamber $c$ in the negative
building has a \emph{codistance} $\delta^*(c)=\delta^*(c_+,c)\in W$,
where $(W,S)$ is the affine Coxeter system. For two chambers in the same
half, their Weyl distance is the element of $W$ represented by the types
of a minimal gallery between them. Write $\ell$ for word length with respect to $S$, and
$W_J=\langle J\rangle$ for $J\subseteq S$.
The chambers of codistance $1$ are exactly the
$B_+$-orbit of the negative base chamber. We use the convention that an
$s$-panel varies the vertex of type $s$. If $\delta^*(c)=w$, the
codistances in this panel belong to $\{w,ws\}$; the value $ws$ is attained,
and, when $\ell(ws)<\ell(w)$, every chamber in the panel other than $c$
has codistance $ws$. These are the twin-building axioms with the positive
argument fixed \cite[Definition~1.1]{AVM2001}.
Every proper $J\subsetneq S$ is spherical:
a $J$-residue, obtained by varying only types in $J$, is a finite spherical
building whose irreducible factors have type $A$ over $k$. In the lattice
model, fixing the other types fixes a partial lattice chain; completing
that chain amounts to choosing flags in its successive finite-dimensional
$k$-vector-space quotients. This gives the asserted type-$A$ factors.

Let $L'$ be the subcomplex consisting of all codistance-$1$ chambers and
their faces. The preceding description shows that $U$ acts simply
transitively on its chambers. We check their vertex stabilizers.
In the standard lattice model at infinity, vertices are homothety classes
of $k[[t^{-1}]]$-lattices. Choose the negative base chamber with vertices
represented by exponent vectors
\[
 (0,\ldots,0)\quad\text{(type $0$)},\qquad
 (\underbrace{0,\ldots,0}_{i},\underbrace{1,\ldots,1}_{N-i})
 \quad\text{(type $i$)}.
\]
For a vertex represented by exponents $(b_a)$, a polynomial matrix in
its stabilizer satisfies
\begin{equation}\label{eq:lattice-bounds}
 \deg g_{ab}\le b_a-b_b.
\end{equation}
This is the condition for preserving the lattice generated by $t^{b_a}e_a$
over $k[[t^{-1}]]$; because the determinant is $1$, preservation of its
homothety class implies preservation of the lattice itself. Intersecting
these bounds with $U$ gives upper unitriangular constant matrices at the
type-$0$ vertex. At type
$i\ne0$ they give two upper unitriangular constant diagonal blocks, cut
after coordinate $i$, a zero upper-right block, and arbitrary entries
in $kt$ in the lower-left block. By \cref{lem:patterns}, these stabilizers
are exactly $H_0,H_i$, respectively. Sending $uH_i$ to the corresponding
vertex identifies
\begin{equation}\label{eq:L-opposite}
 L=\CC(U;(H_i))\cong L'.
\end{equation}
This uses the common-representative definition of faces: a face of $L'$
lies in one of its chambers, and every such chamber is a unique translate
of the base chamber.

\subsection{A filtration by codistance}

We have identified $L$ with an opposition subcomplex. The next step is
to compare it with the whole negative building by adding chambers in
codistance order. The homology of each attachment is controlled by
spherical opposition complexes, for which we use the following input.
In a
spherical building of type $A_r$ over $k$, let $O$ be the subcomplex formed
by chambers opposite a fixed chamber and their faces. For
$q\ge2^{r-1}$,
\begin{equation}\label{eq:spherical-opposition}
 \wtH_j(O;\F)=0\qquad(j<r-1).
\end{equation}
This follows from the integral cone function in
\cite[Corollary~6.10]{OVB}, whose chain contraction identity remains valid
after reduction modulo two. For $r=1$, the required assertion is just
nonemptiness of the opposition complex. For a product of spherical buildings, the
opposition complex is the simplicial join of the opposition complexes
of the factors, and the augmented chain formula for joins gives the
same vanishing below its dimension. We henceforth require $q\ge2^{n-1}$,
which covers every proper residue of our affine building.

\begin{proposition}\label{prop:L-acyclic}
The complex $L$ satisfies \eqref{eq:L-acyclic}.
\end{proposition}

\begin{proof}
Starting with $L$, add all chambers of the negative building in increasing
order of the length of their codistance. Order the finitely many
codistance values of each length arbitrarily. Fix a value $w\ne1$, and put
\[
 D=\{s\in S:\ell(ws)<\ell(w)\},\qquad J=S\setminus D.
\]
The descent set $D$ is nonempty, so $J$ is proper and $W_J$ is finite.
The element $w$ is the minimal representative of $wW_J$; in particular
$\ell(wv)=\ell(w)+\ell(v)$ for $v\in W_J$.

Group the chambers of codistance $w$ by their $J$-residues. Such a residue
is specified by the common face $F$ whose vertex types are $D$. Its
varying part is a spherical building of type $J$.
We~claim that the codistance-$w$ chambers in this residue are precisely
the chambers opposite a single chamber $c_{\max}$ there. Indeed, follow
a reduced expression for the longest element $w_J$ from any chamber of
codistance $w$, choosing at each step the codistance increase permitted
by the panel axiom. This reaches $c_{\max}$ of codistance $ww_J$.
From $c_{\max}$, any minimal gallery within the residue has types forming
a reduced word in $W_J$. Every step decreases codistance length, since
\[
 \ell(ww_Jv)=\ell(w)+\ell(w_J)-\ell(v)\qquad(v\in W_J).
\]
The descent panel axiom thus forces the endpoint at Weyl distance $v$
from $c_{\max}$ to have codistance $ww_Jv$. It has codistance $w$ exactly
when $v=w_J$, which proves the claim. It also proves that no chamber
in this residue has smaller codistance length or a different codistance
of the same length. The same formula proves uniqueness of $c_{\max}$:
a chamber of codistance $ww_J$ has Weyl distance $1$ from it.

Let $O$ denote this opposition complex. The chambers in the residue
now being added, together with their faces, form $F*O$. Its intersection
with the previously added subcomplex is exactly $\partial F*O$.
For one inclusion, a face missing a vertex of type $s\in D$ lies in the
corresponding panel of a codistance-$w$ chamber, and hence in a chamber
of codistance $ws$, which was added earlier. Conversely, a face
containing all of $F$ can lie only in chambers of the same $J$-residue,
so the preceding paragraph shows that it was not present earlier.
Distinct blocks for this $w$ meet only in their attaching loci, since
a face outside such a locus determines its face of types $D$.

The relative chain complex of a block is particularly simple:
\begin{equation}\label{eq:relative-chains}
 C_k(F*O,\partial F*O;\F)
 \cong \wtC_{k-|D|}(O;\F),
 \qquad
 H_k(F*O,\partial F*O;\F)
 \cong \wtH_{k-|D|}(O;\F).
\end{equation}
Only faces containing all vertices of $F$ survive in relative chains;
deleting one of those vertices gives zero, leaving the augmented
boundary in $O$, as illustrated in Figure~\ref{fig:codistance-attachment}.
\begin{figure}[tbp]
\centering
\begin{tikzpicture}[
  every node/.style={font=\small,align=center},
  box/.style={draw,rounded corners=2pt,inner sep=6pt,text width=5.0cm},
  old/.style={box,draw=black!55,fill=black!5},
  new/.style={box,draw=blue!50!black,fill=blue!4},
  arrow/.style={-{Latex[length=2mm]},semithick}
]
\node[box,text width=6.3cm] (block) at (0,2.8)
  {A face $\tau$ of the block $F*O$\\
   $F$ has $|D|$ vertices; $\dim O=|J|-1$};
\node[old] (old) at (-3.3,1.0)
  {$\tau$ omits a vertex of $F$\\
   $\tau\in\partial F*O$: already present};
\node[new] (survives) at (3.3,1.0)
  {$\tau=F\cup\sigma$, with $\sigma\in O$\\
   Survives in relative chains; $\sigma=\varnothing$ is allowed};
\draw[arrow] (block.south west) -- (old.north);
\draw[arrow] (block.south east) -- (survives.north);
\node[old] (deletef) at (-3.3,-1.1)
  {A vertex of $F$ is deleted\\
   The resulting face is zero in relative chains};
\node[new] (deleteo) at (3.3,-1.1)
  {A vertex $v$ of $\sigma$ is deleted\\
   The resulting face is $F\cup(\sigma\setminus\{v\})$};
\draw[arrow] (survives.south west) -- (deletef.north east);
\draw[arrow] (survives.south) -- (deleteo.north);
\node at (0,-2.35)
  {$C_k(F*O,\partial F*O;\F)
    \cong \wtC_{k-|D|}(O;\F)$};
\end{tikzpicture}
\caption{One block of the codistance filtration, shown as a schematic
of its faces and boundary operations. The new relative faces contain
every vertex of $F$. Deleting an $F$-vertex lands in the old subcomplex;
the remaining boundary is the augmented boundary of $O$, shifted by
$|D|$. The complex $O$ is not assumed to be contractible.}
\label{fig:codistance-attachment}
\end{figure}

Since $\dim O=|J|-1$, \eqref{eq:spherical-opposition}
shows that these relative groups vanish for
\[
 k<|D|+|J|-1=|S|-1=n.
\]
If $J=\varnothing$, interpret $O$ as the complex consisting of the empty
face; the pair is $(F,\partial F)$ and the same conclusion holds.

Relative chains for the whole step are the direct sum of the block
relative chains. The long exact sequence of the pair therefore shows
that every step induces an isomorphism in reduced homology in degrees
$j<n-1$. The filtration exhausts the negative affine building, which is
contractible \cite[Theorem~II.10A.4 and Remark~II.10A.6]{BH1999},
and simplicial homology commutes with this increasing union. Thus
$\wtH_j(L;\F)=0$ in precisely the asserted range.
\end{proof}

\begin{corollary}\label{cor:generation}
The simple generators $x_i$, $i\in S$, generate $U$.
\end{corollary}

\begin{proof}
By \cref{prop:L-acyclic}, $L$ is connected. In a connected pure complex,
any two chambers can be joined by a sequence of chambers in which
successive chambers share a vertex. If the chambers labeled $u,v\in U$
share a vertex of type $i$, then $uH_i=vH_i$, so $u^{-1}v\in H_i$.
The chamber action is simply transitive, so following such a sequence
from the base chamber shows that $U=\langle H_i:i\in S\rangle$.
Each $H_i$ is generated by a subset of the $x_j$.
\end{proof}

\section{Homology and the finite quotients}
\label{sec:homology}

We now prove that the finite complexes $T_m$ have no cohomology in
degrees $1,\ldots,d-1$. We write $H_j(G;A)$ for group homology with
trivial coefficients in $A$. The acyclic complex $L$ first compares
the homology of $U$ and its congruence subgroups with simplicial homology.
We then use stabilization to show that the congruence subgroups inherit
the vanishing of $H_j(U;\F)$.

\subsection{An acyclic complex as a partial resolution}
\label{subsec:projective-resolution}

Finite odd-order stabilizers make the simplicial chains projective over
the group algebra. This allows the acyclicity already proved for $L$
to replace a full contractible model in the degrees we need.

\begin{lemma}\label{lem:projective-chains}
Let $h\ge0$ be an integer, and let a group $K$ act on a nonempty
simplicial complex $Z$ without inversions. Suppose every nonempty
simplex stabilizer is finite of odd order.
If $\wtH_j(Z;\F)=0$ for $0\le j\le h$, then
\[
 H_j(K;\F)\cong H_j(C_*(Z;\F)_K)\qquad(0\le j\le h).
\]
Here the subscript denotes coinvariants; this is the chain complex of
simplex orbits, with boundary induced by deleting vertices.
\end{lemma}

\begin{proof}
For every $r\ge0$, the permutation module $C_r(Z;\F)$ is a direct sum
of modules $\F[K/H]$, where $H$ runs over simplex stabilizers. Because
$|H|$ is odd, averaging over $H$ splits the augmentation
$\F[H]\to\F$. Thus $\F$ is a projective $\F[H]$-module, and induction
makes $\F[K/H]$ a projective $\F[K]$-module.
The augmented chains of $Z$ are exact through degree $h$. Retain the
chain modules through degree $h+1$, choose a projective module mapping
onto the kernel of the boundary out of $C_{h+1}(Z;\F)$, and continue
in higher degrees. This gives a
projective resolution of the trivial module $\F$ with the same terms
and differentials in all degrees needed to compute $H_j$, $j\le h$.
Taking coinvariants proves the assertion.
\end{proof}

\begin{proposition}\label{prop:homology-comparisons}
For $0<j<n-1$ and $m\ge2$,
\begin{equation}\label{eq:homology-comparisons}
 H_j(U;\F)=0,\qquad
 H_j(T_m;\F)\cong H_j(\Gamma_N(m);\F).
\end{equation}
Moreover, $T_m$ is connected.
\end{proposition}

\begin{proof}
For a nonempty type set $A\subseteq S$, the simplices of $L$ of type $A$
are the cosets of $\bigcap_{i\in A}H_i=U_{S\setminus A}$. Thus $U$ is
transitive on the faces of each type, and the orbit complex $U\backslash L$
is a single $n$-simplex. Its stabilizers are conjugates of the finite
unitriangular groups in \cref{lem:patterns}, so they have odd order.
Apply \cref{lem:projective-chains} and \cref{prop:L-acyclic} to obtain
the first assertion of \eqref{eq:homology-comparisons}. In the coinvariant
boundary each face is obtained by deleting one vertex, with coefficient
$1$; there are no stabilizer-index factors.

Reduction modulo $t^m$ preserves all the stabilizer intersections by
\cref{lem:patterns}. To check the resulting quotient, let two type-$A$
simplices have representatives $u,v\in U$ and the same reduced vertices.
Then the reduction of $u^{-1}v$ belongs to
$\bigcap_{i\in A}\overline H_i=\overline{U_{S\setminus A}}$.
Choose $h\in U_{S\setminus A}$ with this reduction. The element
$vh^{-1}u^{-1}\in\Gamma_N(m)$ carries the first simplex to the second.
Thus reduction identifies precisely the simplex orbits and gives
\begin{equation}\label{eq:quotient}
 \Gamma_N(m)\backslash L\cong T_m.
\end{equation}
The $\Gamma_N(m)$-stabilizers of nonempty simplices are trivial, since
the reduction is injective on every local group and $\Gamma_N(m)$ is
normal in $U$. The action preserves types. A second application of
\cref{lem:projective-chains} gives the remaining homology comparison.
Connectedness follows from \cref{prop:L-acyclic} and \eqref{eq:quotient}.
\end{proof}

\subsection{From pro-stability to ordinary surjectivity}
\label{subsec:pro-stability}

Our next goal is a surjection on homology at each fixed congruence level.
The available stability theorem concerns inverse systems; odd-index
transfer will supply the additional step from that theorem to an
ordinary surjection.

We first specify its ring conventions. For a commutative nonunital ring
$A$, its \emph{integer unitization} is $\Z\ltimes A=\Z\oplus A$, with
\[
 (u,a)(v,b)=(uv,ub+va+ab),\qquad 1=(1,0).
\]
The projection to $\Z$ defines
\[
 \GL_a(A)=\ker\bigl(\GL_a(\Z\ltimes A)\longrightarrow\GL_a(\Z)\bigr).
\]
A row $b\in(\Z\ltimes A)^\ell$ is \emph{$A$-unimodular} if
$b\equiv e_1\pmod A$ and there is a row $c\equiv e_1\pmod A$ with
$b\cdot c=1$. Here $e_1=(1,0,\ldots,0)$. We use the relative stable
range convention of Vaserstein \cite[Section~2, pp.~384--385]{Vas1969}:
$\sr(A)\le s$ means that, for every such row of length $\ell\ge s+1$,
there are $v_1,\ldots,v_{\ell-1}\in A$ for which
\[
 (b_1+v_1b_\ell,\ldots,b_{\ell-1}+v_{\ell-1}b_\ell)
\]
is $A$-unimodular.

An inverse system $(A_m)$ is \emph{pro-zero} if, for every index, some later
transition to that index is zero. A level map of inverse systems of
abelian groups is a \emph{pro-isomorphism} precisely when its kernel
and cokernel systems are pro-zero. The system $(A_m)$ is \emph{Tor-unital}
if $(\Tor_i^{\Z\ltimes A_m}(\Z,\Z))_m$ is pro-zero for every $i>0$.
Morrow's theorem says that the powers of any ideal in a commutative
Noetherian ring form a Tor-unital system
\cite[Theorem~0.3, arXiv~v2]{Morrow}.
For a Tor-unital inverse system of commutative nonunital rings,
Iwasa's integral homology pro-stability theorem says that, if
$r=\max\{2,\sup_m\sr(A_m)\}<\infty$, then
\begin{equation}\label{eq:iwasa}
 (H_j(\GL_a(A_m);\Z))_m\longrightarrow
 (H_j(\GL_{a+1}(A_m);\Z))_m
 \quad\text{is a pro-isomorphism for }a\ge2j+r-1.
\end{equation}
We cite the published version \cite[Theorem~5.13]{Iwasa}.

\begin{lemma}\label{lem:stabilization}
Let $h\ge0$ be an integer. If $a\ge2h+1$, the inclusion on the first
$a$ coordinates induces
surjections
\[
 H_j(\Gamma_a(m);\F)\longrightarrow H_j(\Gamma_{a+1}(m);\F)
 \qquad(0\le j\le h,\ m\ge1).
\]
\end{lemma}

\begin{proof}
Apply \eqref{eq:iwasa} to $A_m=t^m k[t]$. In this case
$\GL_a(A_m)$ is exactly the group of matrices in $\GL_a(R)$ congruent
to $I$ modulo $t^m$: the entrywise map from the unitization is injective
on $A_m$, and an inverse congruent to $I$ lifts back to the unitization.
Such a matrix has determinant in $R^*=k^*$, congruent to $1$ modulo $t^m$;
hence its determinant is $1$. Thus $\GL_a(A_m)=\Gamma_a(m)$.

The nonunital stable range of every $A_m$ is at most $2$. Here is a direct
verification of the relative row criterion. For a unimodular row
$(b_1,\ldots,b_r)$ with $r\ge3$ and $(b_1,\ldots,b_r)\equiv(1,0,\ldots,0)
\pmod{t^m}$, the polynomial $b_1$ is coprime to $t$. At each irreducible
divisor of $b_1$ at which $b_r$ is nonzero, choose a residue for $v$ avoiding
$-b_2/b_r$. The Chinese remainder theorem allows all these choices with
$v\in t^mR$. Then
$(b_1,b_2+vb_r,b_3,\ldots,b_{r-1})$ is unimodular and still has the same
residue modulo $t^m$: at a divisor where $b_r$ vanishes, original
unimodularity already makes one of $b_2,\ldots,b_{r-1}$ nonzero.
To check the unitization condition, if a row $b$
has $b\cdot c=1$ over $R$ and $b\equiv e_1\pmod{t^m}$, then
\[
 c'=e_1+(1-b_1)c
 \quad\text{satisfies}\quad b\cdot c'=1,
 \qquad c'\equiv e_1\pmod{t^m}.
\]
It therefore lifts to a unimodularity relation over $\Z\ltimes A_m$:
the difference between the lifted dot product and $1$ lies in $A_m$,
where the map to $R$ is injective.
This proves the claimed stable-range bound.

Morrow's theorem and \eqref{eq:iwasa} now give pro-isomorphisms in integral
homology for $0\le j\le h$. The natural universal coefficient sequences
\[
 0\longrightarrow H_j(G;\Z)\otimes\F
 \longrightarrow H_j(G;\F)
 \longrightarrow\Tor_1^{\Z}(H_{j-1}(G;\Z),\F)
 \longrightarrow0
\]
give pro-isomorphisms with $\F$ coefficients in the same range. For
$j\ge1$, both integral degrees $j$ and $j-1$ are covered; degree zero
is immediate. Indeed,
the tensor and Tor functors extend to inverse systems and preserve
isomorphisms in the pro-category, and these levelwise exact sequences
are exact there as well. In particular, the systems of cokernels of the
displayed stabilization maps over $\F$ are pro-zero.

For $m'\ge m$, the index $[\Gamma_a(m):\Gamma_a(m')]$ is finite and odd.
To see this, reduce modulo $t^{m'}$ and filter matrices congruent to $I$
modulo $t^m$ by successive powers of $t$; each successive quotient embeds
into an additive matrix group over $k$. Transfer therefore gives
\[
 H_j(\Gamma_a(m);\F)\xrightarrow{\operatorname{tr}}
 H_j(\Gamma_a(m');\F)\xrightarrow{\operatorname{inc}_*}
 H_j(\Gamma_a(m);\F),
 \qquad \operatorname{inc}_*\operatorname{tr}=1.
\]
Thus all homology transition maps, and consequently all cokernel
transition maps, are surjective. A pro-zero system with surjective
transition maps is zero at every level: a sufficiently late transition
to any chosen term is both zero and surjective. Applying this observation
to the cokernels gives the required ordinary surjections.
\end{proof}

\subsection{Trivial conjugation and vanishing cohomology}
\label{subsec:conjugation}

Stabilization lets us represent every homology class on any chosen set
of $N-2$ coordinates. The two remaining coordinates will make each
elementary conjugation act trivially. This is the finite-level argument
corresponding to the stability-to-conjugation principle in
\cite[Corollary~5.14]{Iwasa}.

\begin{lemma}\label{lem:trivial-action}
For $1\le j\le d-1$, conjugation by $\SL_N(R)$ acts trivially on
$H_j(\Gamma_N(m);\F)$, for every $m\ge1$.
\end{lemma}

\begin{proof}
Our choice of $N$ gives $N-2=2(d-1)+1$. By two applications of
\cref{lem:stabilization}, the subgroup $\Gamma_{N-2}(m)$ on the first
$N-2$ coordinates surjects onto $H_j(\Gamma_N(m);\F)$ in the stated
degrees. Conjugating by permutation matrices gives the same surjectivity
on any specified $N-2$ coordinates, since each such matrix normalizes
$\Gamma_N(m)$ and therefore induces an automorphism of its homology.

Every elementary matrix $e_{ab}(f)$ centralizes the congruence subgroup
on the coordinates complementary to $a,b$. Its conjugation action
therefore fixes the image of a surjection onto $H_j(\Gamma_N(m);\F)$,
so that action is trivial. Elementary matrices generate $\SL_N(R)$
because $R$ is Euclidean, proving the lemma.
\end{proof}

We can now recover all the required congruence-group homology from its
coinvariants under $G_m$. The odd order of $G_m$ makes taking these
coinvariants exact, and the preceding lemma makes their action trivial.

\begin{proposition}\label{prop:quotient-homology}
For every $m\ge2$,
\[
 H_j(\Gamma_N(m);\F)=H_j(T_m;\F)=H^j(T_m;\F)=0
 \qquad(1\le j\le d-1).
\]
\end{proposition}

\begin{proof}
The finite group $G_m=U/\Gamma_N(m)$ has odd order: its quotient modulo
$t$ is the upper unitriangular group over $k$, and the kernel has a
filtration whose factors embed into additive matrix groups over $k$.
Let $P_*$ be a projective $\F[U]$-resolution of the trivial module $\F$.
Its restriction to $\Gamma_N(m)$ is still projective, because $\F[U]$
is free over $\F[\Gamma_N(m)]$. Hence
$(P_*)_{\Gamma_N(m)}$ computes $H_*(\Gamma_N(m);\F)$.
The quotient $G_m$ acts on this chain complex, and its action on homology
is the usual conjugation action. Averaging over the odd-order group
$G_m$ makes its coinvariant functor exact. Since
$P_U=(P_{\Gamma_N(m)})_{G_m}$, taking homology gives
\[
 H_j(U;\F)\cong H_j(\Gamma_N(m);\F)_{G_m}.
\]
By \cref{lem:trivial-action}, the right side equals
$H_j(\Gamma_N(m);\F)$ for $1\le j\le d-1$.
Both homology assertions now follow from
\cref{prop:homology-comparisons}, since $d-1<n-1$.
The cohomology assertion follows by duality over $\F$.
\end{proof}

\section{Coboundary expansion}
\label{sec:expansion}

The cohomology calculation has removed the obstruction to coboundary
expansion. We now obtain a uniform cosystolic inequality from the KMS
construction, verify bounded degree and growth, and compare the weights
on the two skeleta.

\subsection{The quantitative input and its hypotheses}

We recall the Kac--Moody--Steinberg (KMS) construction in the case needed
here. The affine diagram $\widetilde A_n$ has type set $S$, and each
proper subdiagram is a disjoint union of paths of type $A$. Its
\emph{spherical local groups} are the corresponding unitriangular groups
$U_J$, $J\subsetneq S$, with their standard inclusions.
Their \emph{universal completion} is the group obtained by taking the
free product of these groups and identifying each $U_J$ with its image
in $U_K$ whenever $J\subseteq K$. This is the KMS group of type
$\widetilde A_n$ \cite[Definition~4.6]{OVB}.

A finite quotient of the completion defines a \emph{KMS complex} if
its map is injective on every spherical local group and preserves their
pairwise intersections; the complex is formed from the images of the
maximal local groups by the common-intersection coset construction
\cite[Definitions~4.8 and~4.11]{OVB}. A diagram on $n+1$ types is
\emph{$n$-classical} when every proper subdiagram has irreducible
components of classical type $A_r$, $B_r$, $C_r$ or $D_r$
\cite[Definition~4.2]{OVB}. Thus $\widetilde A_n$ is $n$-classical.

We use the following part of the theorem of Oppenheim--Valentiner-Branth
\cite[Theorem~1.2]{OVB}. For fixed $n\ge3$ and sufficiently large $q$,
the $(n-1)$-skeleta $Y$ of finite $n$-dimensional, $n$-classical KMS
complexes over $\mathbb F_q$ satisfy, for some $\eta>0$ independent of
the finite quotient,
\begin{equation}\label{eq:cosystolic}
 \norm{\delta_i f}_{Y,i+1}
 \ge\eta\,\dist_{Y,i}(f,Z^i(Y))
 \qquad(0\le i<n-1).
\end{equation}
The weights are exactly \eqref{eq:weights}. We only need this inequality,
not the additional lower bound on nontrivial cosystoles.

The polynomial local groups above also appear in
\cite[Example~4.13]{OVB}, where the finite quotients use irreducible
polynomials. We verify the recognition conditions for our reductions
modulo $t^m$ directly. Our matrix model supplies
a homomorphism from that completion to $U$, and it is onto by
\cref{cor:generation}. The local groups are realized faithfully by
\cref{lem:patterns}; their intersections are the groups belonging to
the intersecting subsets.

The composite map to $G_m$ has both properties by \cref{lem:patterns}.
Indeed the equality of intersections already holds in the universal
completion: a member of two local groups has, in the matrix model,
image in $U_{J\cap K}$, and injectivity on either local group identifies
it with the corresponding member of that common subgroup.
Therefore $T_m$ is an $n$-dimensional, $n$-classical KMS complex.
This also explains why reduction modulo $t^m$ is allowed: the definition
requires a finite group quotient with the two stated local properties,
and imposes no field condition on the quotient ring.

Choose $q$ once and for all to satisfy $q\ge5$, $q\ge2^{n-1}$, and the
largeness requirement for \eqref{eq:cosystolic}. There are arbitrarily
large odd primes. Thus
\[
 Y_m=\skel_{n-1}T_m
\]
satisfies \eqref{eq:cosystolic} with a common constant $\eta>0$.
It remains to pass from these $(n-1)$-skeleta to the $d$-skeleta $X_m$
without losing a uniform expansion constant.

\subsection{Degree bounds and growth}

Recall $M=q^{\binom N2}$ from \eqref{eq:M}. The common stabilizer of all
base vertices is $U_{\varnothing}=\{1\}$, so the chambers of $T_m$ are
in bijection with $G_m$. A face of type $A$ lies in exactly
$|U_{S\setminus A}|$ chambers. In particular, every vertex lies in exactly
$M$ chambers, and every panel (a face of dimension $n-1$) lies in exactly
$q$ chambers. Each chamber has $N$ vertices, so every vertex of $X_m$
lies in at most
\begin{equation}\label{eq:D}
 D=M\binom{N-1}{d}
\end{equation}
faces of dimension $d$. The complexes $T_m$ and $X_m$ are pure, and
$X_m$ is connected by \cref{prop:homology-comparisons}.

The subgroup $\{e_{12}(f):f\in R\}$ lies in $U$ by definition, and its
image in $G_m$ has order $q^m$. There are $|G_m|/M$ vertices of each
of the $N$ types. Thus
\begin{equation}\label{eq:growth}
 |X_m(0)|=\frac{N|G_m|}{M}\ge\frac{Nq^m}{M}\longrightarrow\infty.
\end{equation}

\subsection{Weights on skeleta}

The expansion theorem uses the intrinsic weights of $Y_m$, while
\cref{thm:main} uses those of $X_m$. The following elementary observation
keeps track of this change of normalization.
Let $T$ be a finite pure $n$-complex, put $X=\skel_s T$, and write
$\mu_i=w_i^T$ and $\nu_i=w_i^X$. For $\tau\in T(s)$ let $c(\tau)$
be its number of containing chambers. If $0<a\le c(\tau)\le b$ for all
$\tau\in T(s)$, then, with
\[
 \Lambda=\frac{\binom{n+1}{s+1}|T(n)|}{|T(s)|},
\]
we have
\begin{equation}\label{eq:skeleton-weights}
 \frac{\Lambda}{b}\mu_i(\sigma)\le\nu_i(\sigma)
 \le\frac{\Lambda}{a}\mu_i(\sigma)
 \qquad(0\le i\le s).
\end{equation}
For $i=s$, this follows from $\nu_s(\tau)=\Lambda\mu_s(\tau)/c(\tau)$.
For $i<s$, use the averaging identities
\[
 \mu_i(\sigma)=\binom{s+1}{i+1}^{-1}
       \sum_{\substack{\tau\in T(s)\\\sigma\subseteq\tau}}\mu_s(\tau),
 \qquad
 \nu_i(\sigma)=\binom{s+1}{i+1}^{-1}
       \sum_{\substack{\tau\in T(s)\\\sigma\subseteq\tau}}\nu_s(\tau).
\]
The first is a double count of chamber--face incidences, and the second
is the definition of skeleton weights. Summing the bounds in degree
$s$ proves \eqref{eq:skeleton-weights}. The same inequalities hold for
norms and distances to any common subspace.

Apply this first with $T=T_m$ and $s=n-1$. Every panel lies in $q$
chambers, so $a=b=\Lambda=q$ and the weights of $Y_m$ are exactly
the weights of $T_m$ in their common degrees. Next take $s=d$.
Every $d$-face lies in between $1$ and $M$ chambers, so
\eqref{eq:skeleton-weights} gives
\begin{equation}\label{eq:norm-comparison}
 \frac{\Lambda_m}{M}\norm{f}_{Y_m,i}
 \le\norm{f}_{X_m,i}
 \le\Lambda_m\norm{f}_{Y_m,i}\qquad(0\le i\le d).
\end{equation}
Here $\Lambda_m=\binom N{d+1}|T_m(n)|/|T_m(d)|$ is common to all degrees
and cancels in the expansion inequality below.

\begin{proof}[Proof of \cref{thm:main}]
Take the complexes in \eqref{eq:construction}. We have proved finiteness,
purity, connectedness, the degree bound \eqref{eq:D}, and the growth
\eqref{eq:growth}. For $i<d$, the cochain spaces and coboundary maps
needed to compute $Z^i$ and $B^i$ are identical for $X_m$, $Y_m$, and
$T_m$. By \cref{prop:quotient-homology}, and by connectedness when $i=0$,
\[
 Z^i(Y_m)=B^i(Y_m)=B^i(X_m)\qquad(0\le i<d).
\]
Combining \eqref{eq:cosystolic} and \eqref{eq:norm-comparison} gives
\begin{align*}
 \norm{\delta_i f}_{X_m,i+1}
 &\ge\frac{\Lambda_m}{M}\norm{\delta_i f}_{Y_m,i+1}\\
 &\ge\frac{\Lambda_m\eta}{M}\dist_{Y_m,i}(f,Z^i(Y_m))\\
 &\ge\frac{\eta}{M}\dist_{X_m,i}(f,B^i(X_m)).
\end{align*}
Thus \eqref{eq:goal} holds with $\eps=\eta/M>0$, independent of $m$
and $i$. Reindexing $m=2,3,\ldots$ gives the required sequence.
\end{proof}

\begin{thebibliography}{99}
\small

\bibitem{Abramenko}
P.~Abramenko,
Group actions on twin buildings,
\emph{Bulletin of the Belgian Mathematical Society -- Simon Stevin}
\textbf{3} (1996), no.~4, 391--406.
\href{https://eudml.org/doc/119921}{EuDML:119921}.

\bibitem{AbramenkoBrown}
P.~Abramenko and K.~S.~Brown,
\emph{Buildings: Theory and Applications},
Graduate Texts in Mathematics \textbf{248}, Springer, 2008.
\href{https://doi.org/10.1007/978-0-387-78835-7}{doi:10.1007/978-0-387-78835-7}.

\bibitem{AVM2001}
P.~Abramenko and H.~Van Maldeghem,
1-twinnings of buildings,
\emph{Mathematische Zeitschrift} \textbf{238} (2001), no.~1, 187--203.
\href{https://doi.org/10.1007/PL00004899}{doi:10.1007/PL00004899}.

\bibitem{BH1999}
M.~R.~Bridson and A.~Haefliger,
\emph{Metric Spaces of Non-Positive Curvature},
Grundlehren der mathematischen Wissenschaften \textbf{319}, Springer, 1999.
\href{https://doi.org/10.1007/978-3-662-12494-9}{doi:10.1007/978-3-662-12494-9}.

\bibitem{CL2025}
M.~Chapman and A.~Lubotzky,
Stability of homomorphisms, coverings and cocycles II: Examples,
applications and open problems,
\emph{Advances in Mathematics} \textbf{463} (2025), article 110117.
\href{https://doi.org/10.1016/j.aim.2025.110117}{doi:10.1016/j.aim.2025.110117}.

\bibitem{DD2024}
Y.~Dikstein and I.~Dinur,
Coboundary and cosystolic expansion without dependence on dimension or degree,
in \emph{APPROX/RANDOM 2024}, LIPIcs \textbf{317}, 62:1--62:24.
\href{https://doi.org/10.4230/LIPIcs.APPROX/RANDOM.2024.62}{doi:10.4230/LIPIcs.\allowbreak APPROX/RANDOM.2024.62}.

\bibitem{EK2024}
S.~Evra and T.~Kaufman,
Bounded degree cosystolic expanders of every dimension,
\emph{Journal of the American Mathematical Society}
\textbf{37} (2024), no.~1, 39--68.
\href{https://doi.org/10.1090/jams/1019}{doi:10.1090/jams/1019}.

\bibitem{GdPVB2025}
L.~Grave de Peralta and I.~Valentiner-Branth,
High-dimensional expanders from Kac--Moody--Steinberg groups,
\emph{European Journal of Combinatorics} \textbf{126} (2025), article 104131.
\href{https://doi.org/10.1016/j.ejc.2025.104131}{doi:10.1016/j.ejc.2025.104131}.

\bibitem{Gromov2010}
M.~Gromov,
Singularities, expanders and topology of maps. Part 2: From combinatorics
to topology via algebraic isoperimetry,
\emph{Geometric and Functional Analysis} \textbf{20} (2010), no.~2, 416--526.
\href{https://doi.org/10.1007/s00039-010-0073-8}{doi:10.1007/s00039-010-0073-8}.

\bibitem{Iwasa}
R.~Iwasa,
Homology pro stability for Tor-unital pro rings,
\emph{Homology, Homotopy and Applications} \textbf{22} (2020), no.~1, 343--374.
\href{https://doi.org/10.4310/HHA.2020.v22.n1.a20}{doi:10.4310/HHA.2020.v22.n1.a20}.

\bibitem{KKL2016}
T.~Kaufman, D.~Kazhdan, and A.~Lubotzky,
Isoperimetric inequalities for Ramanujan complexes and topological expanders,
\emph{Geometric and Functional Analysis} \textbf{26} (2016), no.~1, 250--287.
\href{https://doi.org/10.1007/s00039-016-0362-y}{doi:10.1007/s00039-016-0362-y}.

\bibitem{KM2021}
T.~Kaufman and D.~Mass,
Unique-neighbor-like expansion and group-independent cosystolic expansion,
in \emph{ISAAC 2021}, LIPIcs \textbf{212}, 56:1--56:17.
\href{https://doi.org/10.4230/LIPIcs.ISAAC.2021.56}{doi:10.4230/LIPIcs.ISAAC.2021.56}.

\bibitem{KO2021}
T.~Kaufman and I.~Oppenheim,
Coboundary and cosystolic expansion from strong symmetry,
in \emph{ICALP 2021}, LIPIcs \textbf{198}, 84:1--84:16.
\href{https://doi.org/10.4230/LIPIcs.ICALP.2021.84}{doi:10.4230/LIPIcs.ICALP.2021.84}.

\bibitem{KO2023}
T.~Kaufman and I.~Oppenheim,
High dimensional expanders and coset geometries,
\emph{European Journal of Combinatorics} \textbf{111} (2023), article 103696.
\href{https://doi.org/10.1016/j.ejc.2023.103696}{doi:10.1016/j.ejc.2023.103696}.

\bibitem{KOW2025}
T.~Kaufman, I.~Oppenheim, and S.~Weinberger,
\emph{Coboundary expansion of coset complexes}, 2024;
conference version in \emph{STOC 2025}.
\href{https://arxiv.org/abs/2411.02819v1}{arXiv:2411.02819v1}.

\bibitem{LM2006}
N.~Linial and R.~Meshulam,
Homological connectivity of random 2-complexes,
\emph{Combinatorica} \textbf{26} (2006), no.~4, 475--487.
\href{https://doi.org/10.1007/s00493-006-0027-9}{doi:10.1007/s00493-006-0027-9}.

\bibitem{LLR2019}
A.~Lubotzky, Z.~Luria, and R.~Rosenthal,
Random Steiner systems and bounded degree coboundary expanders of every dimension,
\emph{Discrete \& Computational Geometry} \textbf{62} (2019), no.~4, 813--831.
\href{https://doi.org/10.1007/s00454-018-9991-2}{doi:10.1007/s00454-018-9991-2}.

\bibitem{LM2015}
A.~Lubotzky and R.~Meshulam,
Random Latin squares and 2-dimensional expanders,
\emph{Advances in Mathematics} \textbf{272} (2015), 743--760.
\href{https://doi.org/10.1016/j.aim.2014.12.015}{doi:10.1016/j.aim.2014.12.015}.

\bibitem{LMM2016}
A.~Lubotzky, R.~Meshulam, and S.~Mozes,
Expansion of building-like complexes,
\emph{Groups, Geometry, and Dynamics} \textbf{10} (2016), no.~1, 155--175.
\href{https://doi.org/10.4171/GGD/346}{doi:10.4171/GGD/346}.

\bibitem{MW2009}
R.~Meshulam and N.~Wallach,
Homological connectivity of random $k$-dimensional complexes,
\emph{Random Structures \& Algorithms} \textbf{34} (2009), no.~3, 408--417.
\href{https://doi.org/10.1002/rsa.20238}{doi:10.1002/rsa.20238}.

\bibitem{Morrow}
M.~Morrow,
Pro unitality and pro excision in algebraic $K$-theory and cyclic homology,
\emph{Journal f\"ur die reine und angewandte Mathematik} \textbf{736} (2018), 95--139.
\href{https://arxiv.org/abs/1404.4179v2}{arXiv:1404.4179v2}.

\bibitem{OP2022}
R.~O'Donnell and K.~Pratt,
High-dimensional expanders from Chevalley groups,
in \emph{CCC 2022}, LIPIcs \textbf{234}, 18:1--18:26.
\href{https://doi.org/10.4230/LIPIcs.CCC.2022.18}{doi:10.4230/LIPIcs.CCC.2022.18}.

\bibitem{OVB}
I.~Oppenheim and I.~Valentiner-Branth,
\emph{New cosystolic high-dimensional expanders from KMS groups},
2025, revised September 9, 2026; to appear in \emph{Mathematische Annalen}.
\href{https://arxiv.org/abs/2504.05823v2}{arXiv:2504.05823v2}.

\bibitem{OVBNonAbelian2025}
I.~Oppenheim and I.~Valentiner-Branth,
\emph{Non-Abelian expansion of congruence KMS complexes}, 2025.
\href{https://arxiv.org/abs/2510.18351v1}{arXiv:2510.18351v1}.

\bibitem{Vas1969}
L.~N.~Vaserstein,
On the stabilization of the general linear group over a ring,
\emph{Mathematics of the USSR-Sbornik} \textbf{8} (1969), no.~3, 383--400.
\href{https://doi.org/10.1070/SM1969v008n03ABEH001279}{doi:10.1070/SM1969v008n03ABEH001279}.

\end{thebibliography}

\end{document}
